% Authoritative exercise source. Preserve all derivations and diagrams.

\ExerciseStart{G1.08}{Attraversare i regimi della diffrazione}{Consolidamento}

\BookHeading{Problema, dati e obiettivo}

Per una singola fenditura mantenere \(\lambda=550\,\mathrm{nm}\), \(L=1\,\mathrm m\), e confrontare \(a=20\,\mu\mathrm m\), \(1\,\mathrm{mm}\), \(8\,\mathrm{mm}\). Nel terzo caso usare una media rettangolare del rivelatore di larghezza \(\delta=0{,}16\,\mathrm{mm}\). Determinare il regime e il riferimento geometrico mediato.

\ExTheory{\ExRef{c1-fenditure}{1.04}}{Il numero \(F_a=a^2/(\lambda L)\) confronta apertura e scala di Fresnel \(\sqrt{\lambda L}\). Il campo lontano richiede \(F_a\ll1\). Nel riferimento geometrico la singola apertura trasmette uniformemente per \(|x|<a/2\). Un rivelatore finito ne media il profilo su un intervallo di larghezza \(\delta\).}

\ExPrediction{Al crescere dell’apertura l’interno dell’immagine geometrica si estende. Le oscillazioni ai bordi hanno una scala relativa sempre più piccola; la media del rivelatore ne attenua quelle che non risolve.}

\ExFigure{\begin{tikzpicture}[>=Stealth,exlabel]
\node[exbox] (n0) at (0.0,0.0) {$F_a\ll1$\ campo lontano};\node[exbox] (n1) at (3.6,0.0) {$F_a\simeq1$\ Fresnel};\draw[exother] (n0)--(n1);\node[exbox] (n2) at (7.2,0.0) {$F_a\gg1$\ interno geometrico};\draw[exother] (n1)--(n2);
\end{tikzpicture}}{Il parametro cambia mantenendo fisse lunghezza d’onda e distanza. La propagazione di Fresnel attraversa tutti e tre i casi.}


\begin{ExTasks}

\item Calcolare \(\lambda L\) e la sua radice.

\item Calcolare i tre numeri di Fresnel.

\item Individuare il caso nel quale usare la formula di campo lontano.

\item Costruire la media geometrica vicino al bordo destro \(b=a/2\).

\item Valutare il riferimento in \(b-\delta/2,b,b+\delta/2\) e descrivere il confronto corretto con Fresnel.

\end{ExTasks}

\ExSolution

\ExStep{Determinare la scala}

\ExEq{\lambda L&=5{,}5\cdot10^{-7}\,\mathrm m^2,\\\sqrt{\lambda L}&\simeq7{,}416\cdot10^{-4}\,\mathrm m=0{,}742\,\mathrm{mm}.}
Questa lunghezza caratterizza la struttura diffrattiva trasversale nel modello paraxiale.

\ExStep{Calcolare i tre regimi}

\ExEq{F_{20\mu m}&=\frac{4\cdot10^{-10}}{5{,}5\cdot10^{-7}}\simeq0{,}000727,\\F_{1mm}&=\frac{10^{-6}}{5{,}5\cdot10^{-7}}\simeq1{,}818,\\F_{8mm}&=\frac{64\cdot10^{-6}}{5{,}5\cdot10^{-7}}\simeq116{,}36.}
Le aperture attraversano campo lontano, regime intermedio e regime prossimo alla geometria nell’interno.

\ExStep{Delimitare la formula dei minimi}

Per \(20\,\mu\mathrm m\), \(x_1=\lambda L/a=27{,}5\,\mathrm{mm}\). Inserire \(8\,\mathrm{mm}\) nella stessa espressione darebbe \(0{,}06875\,\mathrm{mm}\), ma \(F_a\gg1\) invalida l’ipotesi utilizzata per ricavarla. Il profilo in quel regime va propagato con Fresnel.

\ExStep{Mediare il riferimento geometrico}

Per \(a=8\,\mathrm{mm}\), il bordo destro è \(b=4\,\mathrm{mm}\). Vicino a quel bordo, la frazione illuminata dell’intervallo del rivelatore è
\ExEq{G_\delta(x)&=\begin{cases}1&x\leq b-\delta/2,\\(b+\delta/2-x)/\delta&b-\delta/2<x<b+\delta/2,\\0&x\geq b+\delta/2.\end{cases}}
La formula locale vale lontano dal bordo sinistro, che è simmetrico.

\ExStep{Valutare e interpretare la media}

Poiché \(\delta/2=0{,}08\,\mathrm{mm}\),
\ExEq{G_\delta(3{,}92)&=1,\\G_\delta(4)&=(4{,}08-4)/0{,}16=1/2,\\G_\delta(4{,}08)&=0.}
Questi sono valori del riferimento geometrico mediato. Il profilo di Fresnel mediato conserva correzioni ai bordi; l’accordo va valutato alla scala dello schermo, senza imporre uguaglianza punto per punto.

\par\noindent\begin{minipage}{\linewidth}
\ExFigure{\begin{tikzpicture}[>=Stealth,exlabel]
\draw[exaxis] (0,0)--(6.2,0) node[right] {$x$ (mm)};\draw[exaxis] (0,0)--(0,2.8) node[above] {$G_\delta$};\draw[Accent,thick] (0,2)--(2,2)--(4,0)--(5.7,0);\draw[Secondary,dashed] (3,0)--(3,1)--(0,1);\node[left] at (0,2) {$1$};\node[left] at (0,1) {$1/2$};\node[below] at (2,0) {$3{,}92$};\node[below] at (3,-.35) {$4$};\node[below] at (4,0) {$4{,}08$};
\end{tikzpicture}}{Ingrandimento del bordo geometrico dopo la stessa media rettangolare del rivelatore; non è un profilo di Fresnel.}


\ExInterpret{Il limite geometrico riguarda un confronto di scale e di osservabili. Una singola apertura larga può ancora diffrangere ai bordi. Per due aperture occorrerebbe controllare anche l’estensione totale illuminata, oltre alla larghezza di ciascuna.}

\ExLab{Nella sezione 1.04 scegliere «Una» e confrontare i preset «Fenditura stretta», «Regime intermedio» e «Limite geometrico». Nell’ultimo controllare \(a=8\,\mathrm{mm}\) e risoluzione \(0{,}16\,\mathrm{mm}\). Confrontare profilo propagato e ombra geometrica dopo la stessa media. La distribuzione dei conteggi è normalizzata sulla finestra osservata.}


\end{minipage}
